\begin{solapartat}
\punts{0,25} Es tracta de crear una força que freni els electrons de manera que no arribin a l'ànode i d'aquesta manera fer que el corrent que circula pel circuit sigui zero. Llavors la força ha de tenir el sentit oposat a la velocitat i com els electrons tenen càrrega negativa el camp elèctric ha de tenir el sentit oposat a la força, tal i com s'indica a la figura:

\figura[0.35]{sol_fig1.png}

Com el camp elèctric apunta en la direcció que disminueix el potencial, el càtode ha d'estar connectat al potencial alt i l'ànode al potencial baix.

\punts{0,5} Quan el potencial aplicat és el potencial de frenada, $V_f$, el treball aplicat pel camp elèctric ha de ser igual a l'energia cinètica màxima dels electrons, així la velocitat dels electrons serà nul·la quan arriben a l'ànode:
\[ E_C = eV_f \]

\punts{0,25} I si apliquem balanç d'energia de l'efecte fotoelèctric tenim:
\[ eV_f = E_C = hf - W_0 \]
On $W_0$ és la funció de treball o treball d'extracció. Llavors:
\[ V_f = \frac{hf - W_0}{e} \]
També es pot expressar com:
\[ f = \frac{eV_f + W_0}{h} \]

\punts{0,25} Podem determinar la constant de Planck a partir de dues parelles de dades de la taula,
\[ \left.\begin{aligned} eV_{f,1} &= hf_1 - W_0\\ eV_{f,2} &= hf_2 - W_0 \end{aligned}\right\} \Rightarrow eV_{f,2} - eV_{f,1} = hf_2 - hf_1 \]
I si aïllem la constant de Planck obtenim
\begin{align*}
h = e\,\frac{V_{f,2} - V_{f,1}}{f_2 - f_1} &= 1,602\times 10^{-19}\,\frac{3,29 - 0,4}{(12,5 - 5,49)\times 10^{14}}\\
&= 6,60\times 10^{-34}\ \mathrm{J\cdot s}
\end{align*}
\end{solapartat}

\begin{solapartat}
\punts{0,75} Un cop determinada la constant de Planck, $W_0$ és:
\[ W_0 = hf - eV_f = 2,98\times 10^{-19}\ \mathrm{J} \]
I la freqüència llindar és
\[ f_0 = \frac{W_0}{h} = 4,52\times 10^{14}\ \mathrm{Hz} \]

\punts{0,25} S'emetran electrons sempre i quan la freqüència de la llum superi a la freqüència llindar.

\punts{0,25} La freqüència de la llum és:
\[ f = \frac{c}{\lambda} = \frac{3,00\times 10^{8}}{350\times 10^{-9}} = 8,57\times 10^{14}\ \mathrm{Hz} \]
Com $f > f_0$ hi haurà emissió d'electrons.

\destaca{Nota pels correctors:} en els criteris generals s'estableix que un error no s'ha de penalitzar dues vegades en el mateix problema. Per tant, no ha de tenir cap mena de penalització el fet que en la resolució d'aquest apartat l'estudiant utilitzi un valor incorrecte de la constant de Planck o que utilitzi el valor conegut de $6,63\times 10^{-34}\ \mathrm{J\cdot s}$.
\end{solapartat}
